\section*{Hoofdstuk \ref{ch:b1:wave-propagation} --- Voortplanting van signalen: lopende en staande golven}

\begin{solution}{exo:b1:wave-propagation:1}
$d = 340 \times 6.0 = \qty{2.0}{km}$. Heen en terug: diepte $= 1500 \times
0.40/2 = \qty{300}{m}$.
\end{solution}

\begin{solution}{exo:b1:wave-propagation:2}
$\lambda = 340/440 = \qty{0.77}{m}$; $\Delta\varphi = 2\pi \times
0.20/0.773 = \qty{1.6}{rad}$ ($93^\circ$); in \omterm{def:b1:wave-propagation:signal}{fase} voor afstanden
$n\lambda = n \times \qty{0.77}{m}$.
\end{solution}

\begin{solution}{exo:b1:wave-propagation:3}
$c = \sqrt{80/4.0\times10^{-3}} = \qty{141}{m/s}$; $\lambda = c/f =
\qty{0.71}{m}$.
\end{solution}

\begin{solution}{exo:b1:wave-propagation:4}
$444 - 440 = 4$ zwevingen per seconde; periode van de omhullende $1/4 = \qty{0.25}{s}$.
Losser draaien verlaagt de snaar: \qty{441}{Hz} (\qty{439}{Hz} bereiken zou
betekenen dat men door de stilte bij \qty{440}{Hz} is gegaan); controleer door
nog wat losser te draaien: de zwevingen moeten verdwijnen en dan terugkeren.
\end{solution}

\begin{solution}{exo:b1:wave-propagation:5}
$f_1 = c/2L = 240/1.2 = \qty{200}{Hz}$; $400$, \qty{600}{Hz}. Derde mode:
knopen op $0$, $0.20$, $0.40$, \qty{0.60}{m}. Raak aan op \qty{20}{cm} (of
\qty{40}{cm}) van een uiteinde: een knoop daar houdt alleen $n = 3, 6, \dots$ over, en
de snaar klinkt op \qty{600}{Hz}.
\end{solution}

\begin{solution}{exo:b1:wave-propagation:6}
$\lambda = \qty{0.40}{m}$. Op het lijnstuk, op afstand $x$ van luidspreker
1, is $\delta = (2.0 - x) - x = 2.0 - 2x$; stilte voor $\delta = \pm\lambda/2,
\pm 3\lambda/2, \dots$: $x = 0.1, 0.3, \dots, 1.9$ m, om de \qty{20}{cm};
het midden is luid. Op de middelloodlijn is $\delta = 0$: overal luid (alleen
de afname met $1/r$). In tegenfase: luid en stil wisselen om --- stilte in het
midden en om de \qty{20}{cm} daarvandaan.
\end{solution}

\begin{solution}{exo:b1:wave-propagation:7}
$i = \lambda D/a = 5.89\times10^{-7} \times 1.5/2.0\times10^{-4} =
\qty{4.4}{mm}$; blauw: \qty{3.4}{mm}. Binnen $\pm\qty{1.0}{cm}$: strepen op
$0$, $\pm 4.4$, $\pm 8.8$ mm --- vijf.
\end{solution}

\begin{solution}{exo:b1:wave-propagation:8}
Gelijk: $I_{\max} = 4I_0$, $I_{\min} = 0$, contrast $1$. $I_0$ en $I_0/4$:
$I_{\max} = I_0(1 + \tfrac14 + 2 \times \tfrac12) = 2.25I_0$, $I_{\min} =
I_0(1.25 - 1) = 0.25I_0$; contrast $2.0/2.5 = 0.8$.
\end{solution}

\begin{solution}{exo:b1:wave-propagation:9}
Gesloten--open: $f = (2n + 1)c/4L = 170$, $510$, \qty{850}{Hz} (alleen oneven
harmonischen: de klarinet). Open--open: $nc/2L = 340$, $680$,
\qty{1020}{Hz} (de fluit).
\end{solution}

\begin{solution}{exo:b1:wave-propagation:10}
$v_\varphi = \sqrt{9.81 \times 10/2\pi} = \qty{3.9}{m/s}$ en
\qty{12.5}{m/s} voor \qty{100}{m}. De lange deining het eerst: $10^6/12.5 =
\qty{22}{h}$ tegen $10^6/3.95 = \qty{70}{h}$, twee dagen eerder. De snelheid
hangt van de golflengte af: een gemengde zee sorteert zichzelf onderweg op
golflengte --- dispersie.
\end{solution}

\begin{solution}{exo:b1:wave-propagation:11}
Mode $n$ heeft de vorm $\sin(n\pi x/L)$. Een afgedwongen knoop in $x = L/3$ vereist
$\sin(n\pi/3) = 0$, dus $n$ een veelvoud van $3$: frequenties $3f_1,
6f_1, \dots$ Tokkelen op $L/2$ geeft de snaar daar een maximale uitwijking,
waaraan mode $n$ alleen kan bijdragen als $\sin(n\pi/2) \neq 0$:
de even modes, die een knoop op $L/2$ hebben, krijgen niets.
\end{solution}

\begin{solution}{exo:b1:wave-propagation:12}
$\lambda = c/f = \qty{3.0}{m}$. Faseverschil $2\pi(2d)/\lambda + \pi$;
destructief wanneer dat gelijk is aan $(2m + 1)\pi$: $2d/\lambda = m$,
$d = m\lambda/2 = 0, 1.5, 3.0, \dots$ m; constructief wanneer $2d/\lambda =
m + \tfrac12$: $d = 0.75, 2.25, \dots$ m. Tunnelwanden en -dak kaatsen
de golf terug: een staand patroon met dode plekken om de halve golflengte langs
de weg van de auto.
\end{solution}

\begin{solution}{pb:b1:wave-propagation:1}
\textbf{1.} $[F/\mu] = M L T^{-2}/(M L^{-1}) = L^2 T^{-2}$: een snelheid in
het kwadraat.

\textbf{2.} Een golf naar $-x$ hangt af van $t + x/c$, dus van
$\omega t + kx$; $k = \omega/c = 2\pi f/c$.

\textbf{3.} $c = \sqrt{0.90/1.0\times10^{-3}} = \qty{30}{m/s}$;
$\lambda = c/f = \qty{0.60}{m}$; $k = 2\pi/\lambda = \qty{10.5}{rad/m}$.

\textbf{4.} $\Delta\varphi = kd = 10.5 \times 0.15 = \qty{1.6}{rad} = \pi/2$.

\textbf{5.} $v_{\max} = A\omega = 2.0\times10^{-3} \times 2\pi \times 50 =
\qty{0.63}{m/s}$, vijftig keer minder dan $c$: de punten van de snaar bewegen
traag terwijl de vorm voortsnelt.

\textbf{6.} Zij loopt naar $+x$: \omterm{def:b1:wave-propagation:signal}{fase} $\omega t - kx$. Aan het vaste
uiteinde is de totale uitwijking op elk moment nul: $s_i(0, t) + s_r(0, t)
= A\cos\omega t + (-A)\cos\omega t = 0$.

\textbf{7.} $\cos(\omega t + kx) - \cos(\omega t - kx) =
-2\sin(\omega t)\sin(kx)$, dus $s = -2A\sin(kx)\sin(\omega t)$.

\textbf{8.} Knopen: $\sin kx = 0$, $x = n\lambda/2$; buiken $x = (n +
\tfrac12)\lambda/2$; naburige knopen $\lambda/2$ uit elkaar.

\textbf{9.} $\sin kL = 0$: $L = n\lambda/2$, dus $f_n = nc/2L$.

\textbf{10.} $n = 2L/\lambda = 2.4/0.60 = 4$: vier buiken, een geheel getal
--- resonantie.

\textbf{11.} $3$ buiken: $\lambda = 2L/3 = \qty{0.80}{m}$, $c = \lambda f
= \qty{40}{m/s}$, $F = \mu c^2 = \qty{1.6}{N}$. $6$ buiken: $\lambda =
\qty{0.40}{m}$, $c = \qty{20}{m/s}$, $F = \qty{0.40}{N}$.

\textbf{12.} $c = 2Lf/n = 2 \times 1.20 \times 50/3 = \qty{40}{m/s}$,
$\mu = F/c^2 = 1.6/1600 = \qty{1.0e-3}{kg/m} = \qty{1.0}{g/m}$.

\textbf{13.} $\mu = mgn^2/(4L^2f^2)$: $u(\mu)/\mu = \sqrt{(u_m/m)^2 +
(2u_L/L)^2} = \sqrt{0.61^2 + 0.33^2}\,\% = 0.7\%$; de massa is bepalend.

\textbf{14.} De grondtoon. $c = 2Lf_1$: lage E $2 \times 0.650 \times
82.4 = \qty{107}{m/s}$; hoge E \qty{428}{m/s}.

\textbf{15.} $F = \mu c^2$: lage E $6.0\times10^{-3} \times 107^2 =
\qty{69}{N}$; hoge E $4.0\times10^{-4} \times 428^2 = \qty{73}{N}$ ---
vergelijkbare spankrachten. $f$ met een factor $4$ verlagen bij vaste $L$ en $\mu$
zou $16$ keer minder spankracht vragen: een slappe, rammelende snaar. Omwinden
verhoogt in plaats daarvan $\mu$.

\textbf{16.} $f_n = 2^{n/12}f_1$ en $f \propto 1/L$: $L_n = 2^{-n/12}L$,
afstand vanaf de topkam $L - L_n = L(1 - 2^{-n/12})$: fret 1 \qty{36.5}{mm},
fret 5 \qty{163}{mm}, fret 7 \qty{216}{mm}, fret 12 \qty{325}{mm}.

\textbf{17.} Elke afstand is $2^{-1/12} = 0.944$ keer de vorige.

\textbf{18.} Factor $2^{3/12} = 1.19$ (drie halve tonen); lengte
$650 \times 2^{-3/12} = \qty{547}{mm}$.

\textbf{19.} Aanraken op $L/2$: alleen even modes, toonhoogte $2f_1 =
\qty{165}{Hz}$; op $L/3$: veelvouden van $3$, toonhoogte $3f_1 = \qty{247}{Hz}$.

\textbf{20.} Vlak bij een uiteinde groeit $\sin(n\pi x/L) \approx n\pi x/L$ met
$n$: hoge modes worden verhoudingsgewijs sterker aangeslagen --- helder. Boven het
klankgat (dicht bij het midden) zijn de even modes zwak en overheerst de grondtoon
--- zacht.

\textbf{21.} $3f_1 = \qty{330}{Hz}$ zweeft met \qty{2}{Hz} tegen $2f_2$:
$f_2 = 166$ of \qty{164}{Hz}, verhouding $1.509$ of $1.491$ (de gelijkzwevend
gestemde kwint is $2^{7/12} = 1.498$).

\textbf{22.} De vierde harmonische, $4f$. Drie zwevingen: $4f = 437$ of
\qty{443}{Hz}, $f = 109.25$ of \qty{110.75}{Hz}.

\textbf{23.} De zwevingen versnellen bij aanspannen, dus stond de snaar te hoog:
\qty{110.75}{Hz}. Zij moet $0.75/110.75 = 0.68\%$ dalen, dus de
spankracht $1.4\%$.

\textbf{24.} $f = (n/2L)\sqrt{F/\mu}$, dus $\ln f = \tfrac12\ln F + \text{const}$
en $\dd f/f = \tfrac12\,\dd F/F$.

\textbf{25.} $f_1 = \dfrac{1}{2L}\sqrt{\dfrac{F}{\mu}}$: de lengte legt
$\lambda_1 = 2L$ vast, spankracht en massa leggen $c$ vast, de fretten schalen $L$ met
$2^{-n/12}$, en zwevingen onthullen $f - f_{\text{ref}}$. Achter dat alles staat één
betrekking: $\lambda f = c$, met de staandegolfvoorwaarde $L = n\lambda/2$.
\end{solution}
